\chapter*{Appendix}
\addcontentsline{toc}{chapter}{Appendix}

\section*{A. Prompts for Uncertainty Estimation (Study 1)}

The following prompts were used for complexity estimation in the MMLU-Pro evaluation pipeline.

\subsection*{A.1 Single-Token Response Prompts}

\textbf{Standard prompt:}
\begin{quote}
\textit{The following are multiple choice questions about subject.\\
Write down ONLY the NUMBER of the correct answer and nothing else.}
\end{quote}

\textbf{With fallback for unknown answers:}
\begin{quote}
\textit{The following are multiple choice questions about subject.\\
If you are certain about the answer return the correct option number, otherwise return 0. Write down ONLY the NUMBER and nothing else.}
\end{quote}

\subsection*{A.2 Chain-of-Thought Response Prompts}

\textbf{Standard CoT prompt:}
\begin{quote}
\textit{The following are multiple choice questions about subject. Explain your thinking process step-by-step. At the end, write down the number of the correct answer by strictly following this format: [number of correct answer].}
\end{quote}

\textbf{CoT with fallback:}
\begin{quote}
\textit{The following are multiple choice questions about subject. Explain your thinking process step-by-step. At the end, if you are certain about the answer write down the number of the correct answer by strictly following this format: [number of correct answer], otherwise return [0].}
\end{quote}

\subsection*{A.3 MASJ Prompts}

\textbf{Reasoning requirement estimation:}
\begin{quote}
\textit{You are an expert in the topic of the question. Please act as an impartial judge and evaluate the complexity of the multiple-choice question with options below. Begin your evaluation by providing a short explanation. Be as objective as possible. After providing your explanation, you must rate the question complexity by strictly following the criteria: [Number of reasoning steps] -- how many reasoning steps do you need to answer this question? Valid answers: low, medium, high.}
\end{quote}

\textbf{Education level estimation:}
\begin{quote}
\textit{You are an expert in the topic of the question. Please act as an impartial judge and evaluate the complexity of the multiple-choice question with options below. After providing your explanation, you must rate the question complexity by strictly following the scale: ``high school and easier'', ``undergraduate'', ``graduate'', ``postgraduate''.}
\end{quote}

\section*{B. Training Hyperparameters (Study 3)}

\begin{table}[h]
\centering
\caption{Training hyperparameters for MMLU-Pro experiments.}
\begin{tabular}{lcccc}
\toprule
Parameter & SFT & Pipeline (easy/mid) & Pipeline (hard) & Distillation \\
\midrule
Effective batch size & 64 & 64 & 64 & 64 \\
Max new tokens & 30 & 30 & 8192 & 8192 \\
Optimizer & AdamW & AdamW & AdamW & AdamW \\
Learning Rate & 1e-4 & 1e-4 & 1e-4 & 1e-4 \\
\bottomrule
\end{tabular}
\label{tab:hyperparams_mmlu}
\end{table}

\begin{table}[h]
\centering
\caption{Training hyperparameters for GSM8K/MedMCQA experiments.}
\begin{tabular}{lcccc}
\toprule
Parameter & SFT & Pipeline (easy/mid) & Pipeline (hard) & Distillation \\
\midrule
Effective batch size & 256 & 256 & 256 & 256 \\
Max new tokens & 100 & 100 & 8192 & 8192 \\
Optimizer & AdamW & AdamW & AdamW & AdamW \\
Learning Rate & 1e-4 & 1e-4 & 1e-4 & 1e-4 \\
LoRA rank & 64 & 64 & 64 & 64 \\
LoRA alpha & 128 & 128 & 128 & 128 \\
LoRA dropout & 0.05 & 0.05 & 0.05 & 0.05 \\
\bottomrule
\end{tabular}
\label{tab:hyperparams_gsm8k}
\end{table}

\section*{C. Steering Coefficients (Study 2)}

\begin{table}[h]
\centering
\caption{Llama-3.2-1B steering coefficients: KL-optimized (from distributional analysis) vs.\ CSI-optimized (from generation grid search).}
\begin{tabular}{lcc}
\toprule
Language pair & KL-opt. & CSI-opt. \\
\midrule
English -- Chinese  & $-5.0$ & $-5.0$ \\
English -- Spanish  & $-3.7$ & $-5.0$ \\
English -- Russian  & $-3.5$ & $-4.0$ \\
English -- Hindi    & $-3.6$ & $-5.0$ \\
\bottomrule
\end{tabular}
\label{tab:appendix_steering_coeff}
\end{table}

CSI-optimized coefficients tend toward higher magnitudes, which improves language control but increases degeneration risk. The coefficients cluster near the maximum tested value ($-5.0$), suggesting that direct optimization of generation metrics pushes toward an aggressive regime where language control and fluency trade off sharply.

\section*{D. Supplementary Distributional Metrics (Study 2)}

\begin{table}[h]
\centering
\caption{JS divergence (no steering $\rightarrow$ steering, lower is better).}
\begin{tabular}{lcc}
\toprule
Language pair & Qwen & Llama \\
\midrule
English -- Chinese & 0.67 $\rightarrow$ 0.49 & 0.65 $\rightarrow$ 0.51 \\
English -- Spanish & 0.59 $\rightarrow$ 0.52 & 0.58 $\rightarrow$ 0.44 \\
English -- Russian & 0.62 $\rightarrow$ 0.50 & 0.62 $\rightarrow$ 0.43 \\
English -- Hindi & 0.64 $\rightarrow$ 0.60 & 0.63 $\rightarrow$ 0.54 \\
\bottomrule
\end{tabular}
\label{tab:appendix_js}
\end{table}

\begin{table}[h]
\centering
\caption{Top-50 token overlap (no steering $\rightarrow$ steering, higher is better).}
\begin{tabular}{lcc}
\toprule
Language pair & Qwen & Llama \\
\midrule
English -- Chinese & 3.69\% $\rightarrow$ 26.46\% & 4.47\% $\rightarrow$ 24.31\% \\
English -- Spanish & 11.41\% $\rightarrow$ 20.23\% & 11.45\% $\rightarrow$ 30.74\% \\
English -- Russian & 7.41\% $\rightarrow$ 25.45\% & 7.49\% $\rightarrow$ 31.96\% \\
English -- Hindi & 11.23\% $\rightarrow$ 11.68\% & 6.54\% $\rightarrow$ 25.88\% \\
\bottomrule
\end{tabular}
\label{tab:appendix_overlap}
\end{table}
